% GENERATED FILE. Edit canonical lessons, then run npm run generate:books.
% canonical-source-hash: fnv1a64:eb361e860ef8f9bb
% canonical-lessons: FR-C323-regner, FR-C323-regne, FR-C323-trone, FR-C323-couronne, FR-C323-couronner, FR-R323-revision-316-323

\chapter{Le règne}
\label{ch:fr-le-regne}
\hlchaptermodality{323}

\begin{chapteropening}
\textbf{By the end of this chapter:} \emph{I can tell the story of a king or queen: reigning, the reign, the throne, the crown and crowning.}

Complete the last lesson of chapter 323: I can tell the story of a king or queen: reigning, the reign, the throne, the crown and crowning.
\end{chapteropening}

\section[régner]{régner — to reign, to rule}

\label{lesson:FR-C323-regner}



\begin{quote}
Before the new one: say the French for the myth, then the French for the moral.
\end{quote}

\subsection*{You'll want to know: régner}
\textbf{régner} — \textquotedblleft{}to reign, to rule\textquotedblright{}.

\begin{quote}
Ce roi a régné pendant quarante ans.
\end{quote}

\begin{quote}
\emph{This king reigned for forty years.}
\end{quote}

Figurative: \emph{le silence régnait} — silence reigned.

\begin{cousinweb}
\emph{régner} comes from Old French \emph{reignier}, from Latin \emph{rēgnāre} \textquotedblleft{}to rule\textquotedblright{} (EtymDB), built on \emph{rēgnum} \textquotedblleft{}kingship\textquotedblright{}, from \emph{rēx} \textquotedblleft{}king\textquotedblright{} (Lewis \& Short).
\end{cousinweb}

\begin{grammarlens}[title={What you've built}]
One more word for telling, explaining and describing at B1.
\end{grammarlens}

\subsection*{Your turn}
\begin{itemize}
\raggedright
  \item \emph{Say it:} \emph{régner}
  \item \emph{Say it:} \emph{régner}, once more
  \item \emph{Say it:} a sentence of your own with \emph{régner}
  \item \emph{Recall:} say the French for the myth, then the French for the moral, then say \emph{régner} again
\end{itemize}

\begin{tcolorbox}[breakable,colback=teal!4,colframe=teal!35!black,title={Before you move on}]
What does régner mean? (\textquotedblleft{}To reign, to rule\textquotedblright{}.) Say it once more.
\end{tcolorbox}
\section[le règne]{le règne — the reign}

\label{lesson:FR-C323-regne}



\begin{quote}
Before the new one: say the French for the moral, then the French for to reign.
\end{quote}

\subsection*{You'll want to know: le règne}
\textbf{le règne} — \textquotedblleft{}the reign\textquotedblright{}.

\begin{quote}
Pendant son règne, le pays a connu la paix.
\end{quote}

\begin{quote}
\emph{During his reign, the country knew peace.}
\end{quote}

Also a kingdom in nature: \emph{le règne animal}.

\begin{cousinweb}
\emph{règne} comes from Old French \emph{regne}, from Latin \emph{rēgnum} \textquotedblleft{}kingship, authority\textquotedblright{} (EtymDB).
\end{cousinweb}

\begin{grammarlens}[title={What you've built}]
One more word for telling, explaining and describing at B1.
\end{grammarlens}

\subsection*{Your turn}
\begin{itemize}
\raggedright
  \item \emph{Say it:} \emph{le règne}
  \item \emph{Say it:} \emph{le règne}, once more
  \item \emph{Say it:} a sentence of your own with \emph{le règne}
  \item \emph{Recall:} say the French for the moral, then the French for to reign, then say \emph{le règne} again
\end{itemize}

\begin{tcolorbox}[breakable,colback=teal!4,colframe=teal!35!black,title={Before you move on}]
What does le règne mean? (\textquotedblleft{}The reign\textquotedblright{}.) Say it once more.
\end{tcolorbox}
\section[le trône]{le trône — the throne}

\label{lesson:FR-C323-trone}



\begin{quote}
Before the new one: say the French for to reign, then the French for the reign.
\end{quote}

\subsection*{You'll want to know: le trône}
\textbf{le trône} — \textquotedblleft{}the throne\textquotedblright{}.

\begin{quote}
Le fils du roi est monté sur le trône.
\end{quote}

\begin{quote}
\emph{The king's son came to the throne.}
\end{quote}

\emph{Monter sur le trône} — to come to the throne.

\begin{cousinweb}
\emph{trône} comes from Old French \emph{trone}, from Latin \emph{thronus} \textquotedblleft{}throne\textquotedblright{} (EtymDB), a borrowing from Greek \emph{thrónos} \textquotedblleft{}a seat\textquotedblright{}.
\end{cousinweb}

\begin{grammarlens}[title={What you've built}]
One more word for telling, explaining and describing at B1.
\end{grammarlens}

\subsection*{Your turn}
\begin{itemize}
\raggedright
  \item \emph{Say it:} \emph{le trône}
  \item \emph{Say it:} \emph{le trône}, once more
  \item \emph{Say it:} a sentence of your own with \emph{le trône}
  \item \emph{Recall:} say the French for to reign, then the French for the reign, then say \emph{le trône} again
\end{itemize}

\begin{tcolorbox}[breakable,colback=teal!4,colframe=teal!35!black,title={Before you move on}]
What does le trône mean? (\textquotedblleft{}The throne\textquotedblright{}.) Say it once more.
\end{tcolorbox}
\section[la couronne]{la couronne — the crown; the wreath}

\label{lesson:FR-C323-couronne}



\begin{quote}
Before the new one: say the French for the reign, then the French for the throne.
\end{quote}

\subsection*{You'll want to know: la couronne}
\textbf{la couronne} — \textquotedblleft{}the crown; the wreath\textquotedblright{}.

\begin{quote}
Le roi portait une couronne d'or.
\end{quote}

\begin{quote}
\emph{The king wore a gold crown.}
\end{quote}

Also a ring of flowers or of bread: \emph{une couronne de fleurs}.

\begin{cousinweb}
\emph{couronne} is inherited from Latin \emph{corōna} \textquotedblleft{}a wreath, a crown\textquotedblright{} (EtymDB), itself a borrowing from Greek.
\end{cousinweb}

\begin{grammarlens}[title={What you've built}]
One more word for telling, explaining and describing at B1.
\end{grammarlens}

\subsection*{Your turn}
\begin{itemize}
\raggedright
  \item \emph{Say it:} \emph{la couronne}
  \item \emph{Say it:} \emph{la couronne}, once more
  \item \emph{Say it:} a sentence of your own with \emph{la couronne}
  \item \emph{Recall:} say the French for the reign, then the French for the throne, then say \emph{la couronne} again
\end{itemize}

\begin{tcolorbox}[breakable,colback=teal!4,colframe=teal!35!black,title={Before you move on}]
What does la couronne mean? (\textquotedblleft{}The crown; the wreath\textquotedblright{}.) Say it once more.
\end{tcolorbox}
\section[couronner]{couronner — to crown}

\label{lesson:FR-C323-couronner}



\begin{quote}
Before the new one: say the French for the throne, then the French for the crown.
\end{quote}

\subsection*{You'll want to know: couronner}
\textbf{couronner} — \textquotedblleft{}to crown\textquotedblright{}.

\begin{quote}
Le nouveau roi a été couronné dans la grande église.
\end{quote}

\begin{quote}
\emph{The new king was crowned in the great church.}
\end{quote}

Figurative: \emph{couronner une carrière} — to crown a career.

\begin{cousinweb}
\emph{couronner} is inherited from Latin \emph{corōnāre} \textquotedblleft{}to crown\textquotedblright{} (EtymDB), built on \emph{corōna}.
\end{cousinweb}

\begin{grammarlens}[title={What you've built}]
One more word for telling, explaining and describing at B1.
\end{grammarlens}

\subsection*{Your turn}
\begin{itemize}
\raggedright
  \item \emph{Say it:} \emph{couronner}
  \item \emph{Say it:} \emph{couronner}, once more
  \item \emph{Say it:} a sentence of your own with \emph{couronner}
  \item \emph{Recall:} say the French for the throne, then the French for the crown, then say \emph{couronner} again
\end{itemize}

\begin{tcolorbox}[breakable,colback=teal!4,colframe=teal!35!black,title={Before you move on}]
What does couronner mean? (\textquotedblleft{}To crown\textquotedblright{}.) Say it once more.
\end{tcolorbox}
\section[(dialogue)]{Retour sur les chapitres 316–323}

\label{lesson:FR-R323-revision-316-323}



\begin{quote}
Say the French for to celebrate and the celebration.
\end{quote}

\subsection*{Your turn}
Say these aloud:

\begin{itemize}
\raggedright
  \item to celebrate, the celebration, the procession, the brass band, the public holiday
  \item the parade, to march past, the carnival, the float, to dress up
  \item the fireworks, the firecracker, the dance, the garland, to decorate
  \item the festive late-night dinner, to have the festive late-night dinner, the fir tree, the log, the king cake
  \item to clink glasses, the toast, to wish, the candle, the wish
  \item once upon a time, the fairy, the witch, the dragon, the fable
  \item the account, to narrate, the legend, the myth, the moral
  \item to reign, the reign, the throne, the crown, to crown
\end{itemize}

\begin{tcolorbox}[breakable,colback=teal!4,colframe=teal!35!black,title={Before you move on}]
To celebrate? (\textbf{célébrer}.) The celebration? (\textbf{la célébration}.) The procession? (\textbf{le cortège}.) The brass band? (\textbf{la fanfare}.) The public holiday? (\textbf{le jour férié}.) The parade? (\textbf{le défilé}.) To march past? (\textbf{défiler}.) The carnival? (\textbf{le carnaval}.) The float? (\textbf{le char}.) To dress up? (\textbf{déguiser}.) The fireworks? (\textbf{le feu d'artifice}.) The firecracker? (\textbf{le pétard}.) The dance? (\textbf{le bal}.) The garland? (\textbf{la guirlande}.) To decorate? (\textbf{décorer}.) The festive late-night dinner? (\textbf{le réveillon}.) To have the festive late-night dinner? (\textbf{réveillonner}.) The fir tree? (\textbf{le sapin}.) The log? (\textbf{la bûche}.) The king cake? (\textbf{la galette}.) To clink glasses? (\textbf{trinquer}.) The toast? (\textbf{le toast}.) To wish? (\textbf{souhaiter}.) The candle? (\textbf{la bougie}.) The wish? (\textbf{le vœu}.) Once upon a time? (\textbf{il était une fois}.) The fairy? (\textbf{la fée}.) The witch? (\textbf{la sorcière}.) The dragon? (\textbf{le dragon}.) The fable? (\textbf{la fable}.) The account? (\textbf{le récit}.) To narrate? (\textbf{narrer}.) The legend? (\textbf{la légende}.) The myth? (\textbf{le mythe}.) The moral? (\textbf{la morale}.) To reign? (\textbf{régner}.) The reign? (\textbf{le règne}.) The throne? (\textbf{le trône}.) The crown? (\textbf{la couronne}.) To crown? (\textbf{couronner}.)
\end{tcolorbox}
